% Rebuttal to Reviewer PtZf
% ACM-BCB 2026 Submission: Structure-Aware Prediction of PROTAC-Mediated
% Protein Degradability via Graph Neural Networks
%
% Compile: pdflatex rebuttal_PtZf.tex && pdflatex rebuttal_PtZf.tex

\documentclass[11pt]{article}
\usepackage[margin=1in]{geometry}
\usepackage{xcolor}
\usepackage{booktabs}
\usepackage{enumitem}
\usepackage{hyperref}
\usepackage{amsmath}

% Styling
\definecolor{revblue}{RGB}{0,70,170}
\newcommand{\reviewerbox}[2]{%
  \noindent\fcolorbox{black}{gray!10}{\parbox{\dimexpr\linewidth-2\fboxsep-2\fboxrule}{%
    \textbf{#1}: #2}}%
  \vspace{4pt}%
}
\newcommand{\resp}{\noindent\textbf{Response:}~}
\newcommand{\change}[1]{\textcolor{revblue}{[Paper reference: #1]}}

\setlength{\parskip}{6pt}
\setlength{\parindent}{0pt}

\begin{document}

\begin{center}
{\Large\bfseries Response to Reviewer PtZf}\\[6pt]
{\large Structure-Aware Prediction of PROTAC-Mediated Protein Degradability\\via Graph Neural Networks}\\[6pt]
ACM-BCB 2026 Submission\\[4pt]
\textit{Rating: 3 --- Borderline Reject}
\end{center}

\bigskip

We thank Reviewer PtZf for the detailed feedback. We have addressed all concerns and believe the revised manuscript substantially strengthens the empirical evaluation. All changes are \textcolor{revblue}{highlighted in blue} in the revised paper. Line numbers below refer to the revised \texttt{paper.tex}.

\bigskip
\hrule
\bigskip

%% ====================================================================
%%  W1
%% ====================================================================

\reviewerbox{W1}{Performance is inconsistent across the paper: abstract reports 0.78 AUROC (95\% CI 0.74--0.82), main results section reports 0.646$\pm$0.124, appendix reports 5-fold CV mean 0.565.}

\resp We believe the ``0.78'' figure may refer to a prior draft; the current abstract (line 65) reports 0.646$\pm$0.124 (multi-seed average). To prevent any confusion, we have added a ``Reconciling performance estimates'' paragraph that explicitly maps each reported number to its evaluation protocol:

\begin{center}
\begin{tabular}{llll}
\toprule
\textbf{Number} & \textbf{Protocol} & \textbf{Role} & \textbf{Line} \\
\midrule
0.603 $\pm$ 0.097 & 6-seed, single split & Most conservative & L595 \\
0.646 $\pm$ 0.124 & 3-seed, single split & Primary result & L440 \\
0.7449 & Best seed (seed 42) & Peak (secondary) & L442 \\
0.657 & Baseline (seed 42) & Single-seed reference & L433 \\
0.565 $\pm$ 0.052 & 5-fold CV & Conservative lower bound & L651 \\
\bottomrule
\end{tabular}
\end{center}

The 6-seed mean (0.603) is recommended for citation as the most conservative estimate. The differences between protocols are now explicitly reconciled:

\begin{quote}\itshape
``We report three evaluation protocols: multi-seed single split (0.646$\pm$0.124), baseline single seed (0.657), and 5-fold CV (0.565$\pm$0.052, Appendix Table~9). The 6-seed mean (0.603) is the most conservative and recommended estimate.''
\end{quote}

\change{Lines 650--654 (``Reconciling performance estimates'' paragraph in \S4.3); Abstract line 65 (primary number clearly stated); Table~1 lines 421--438 (all protocols in one table).}

\bigskip

%% ====================================================================
%%  W2
%% ====================================================================

\reviewerbox{W2}{Robustness is limited. Target-unseen has high seed sensitivity, and 5-fold CV mean is substantially lower than the headline single-split number.}

\resp We agree and have substantially expanded the variance analysis:

\begin{enumerate}[topsep=2pt,itemsep=1pt]
  \item \textbf{Extended to 6 seeds} (789, 1011, 1213 added), yielding 0.603$\pm$0.097---marginally below gradient boosting (0.607).
  \change{Lines 593--599 (``Six-seed extended validation'' paragraph in \S4.3).}

  \item \textbf{Stated explicitly that the improvement is not statistically significant}: ``bootstrap p = 0.259'' for 3-seed (line 507) and ``p = 0.556'' for 6-seed (line 595). Previously the $p$-value appeared only in a figure caption; it is now in the main text.
  \change{Lines 441 (``not statistically significant at $\alpha{=}0.05$''), 507 (``bootstrap p = 0.259''), 595 (``bootstrap p = 0.556'').}

  \item \textbf{Recommended ensembling} across $\geq$6 seeds for reliable deployment, noting individual seeds span a 0.25-AUROC range (0.502--0.745).
  \change{Line 598 (``model ensembling across $\geq$6 seeds is strongly recommended'').}

  \item \textbf{Reported both 3-seed and 6-seed} results in Table~1 to show how the estimate changes with more seeds.
  \change{Table~1 lines 421--438 (rows for 3-seed and 6-seed).}

  \item \textbf{Figure~4} (line 605) now visualizes all 6 seeds, distinguishing original (solid) vs.\ additional (hatched).
  \change{Figure~4 at line 605.}

  \item \textbf{Honest performance assessment} in Discussion: ``We emphasize that the 6-seed mean (0.603) is not statistically superior to gradient boosting (0.607, p = 0.556).''
  \change{Lines 713--717 (``Honest performance assessment'' paragraph in \S5).}
\end{enumerate}

\bigskip

%% ====================================================================
%%  W3
%% ====================================================================

\reviewerbox{W3}{Validation and model-selection protocol appears weak for the main use case. Validation AUROC does not reliably predict target-unseen test AUROC.}

\resp Acknowledged as a methodological limitation. We now discuss this explicitly in the Training Procedure section:

\begin{quote}\itshape
``We observe that validation-split AUROC does not reliably predict target-unseen test AUROC. Ideally, we would use a target-unseen validation split for hyperparameter selection; however, this would further reduce the already-limited training data (155 proteins $\to$ ${\sim}$120 training, ${\sim}$15 validation, ${\sim}$20 test). We therefore use standard random validation\ldots while acknowledging this may bias hyperparameter selection toward random-split performance.''
\end{quote}

We also state the implication: the chosen hyperparameters (LR=$5\times10^{-4}$, dropout=0.05, noted at line 310) may not be optimal for target-unseen generalization, contributing to variance.

\change{Lines 314--318 (``Validation-test mismatch'' paragraph in \S3.5 Training Procedure); Line 310 (hyperparameter values); Limitation item~4 at line 735 (``Validation-test mismatch: Validation AUROC does not reliably predict target-unseen test AUROC'').}

\bigskip

%% ====================================================================
%%  W4
%% ====================================================================

\reviewerbox{W4}{Generalization across E3 ligases remains limited. VHL performance is poor; rare ligases have too few samples.}

\resp Fully addressed. E3 claims are narrowed to CRBN$\leftrightarrow$VHL transfer throughout. Key additions:

\begin{itemize}[topsep=2pt,itemsep=1pt]
  \item \textbf{Per-E3 Table~5} (lines 622--637): Quantitative breakdown---CRBN 0.758 ($n{=}273$), VHL 0.396 ($n{=}187$), cIAP1 0.417 ($n{=}8$).
  \change{Table~5 at line 625.}

  \item \textbf{Per-E3 Figure~5} (lines 639--645): Visual bar chart with random-chance reference line.
  \change{Figure~5 at line 644.}

  \item \textbf{VHL root cause analysis} (lines 612--620): Ruled out class imbalance (VHL 36.9\% positive vs CRBN 38.4\%) and dataset size (91 vs 125 unique proteins). Identified structural heterogeneity---VHL-targeted proteins span diverse protein families with varying surface topology.
  \change{Lines 612--620.}

  \item \textbf{Explicit practitioner warning} (line 619): ``practitioners should expect substantially lower performance than the overall 0.646 mean when deploying on VHL-targeted proteins, and near-random performance for rare E3s.''
  \change{Line 619.}

  \item \textbf{E3-unseen vs target-unseen paradox explained} (lines 608--611): VHL can succeed at E3-level transfer (0.811) while failing at protein-level generalization (0.396)---different generalization axes.
  \change{Lines 608--611.}

  \item \textbf{Narrowed throughout}: Abstract line 65 (``CRBN$\to$VHL''), section heading line 608 (``limited to major ligases''), conclusion line 749.
  \change{Lines 65, 608, 749.}
\end{itemize}

\bigskip

%% ====================================================================
%%  W5
%% ====================================================================

\reviewerbox{W5}{No external validation reduces confidence in real-world generalization.}

\resp Stated as Limitation item~6:

\begin{quote}\itshape
``No external validation: PROTAC-8K is the only publicly available benchmark with sufficient scale; generalization to data from other assay platforms or research groups is unvalidated.''
\end{quote}

We also caution against deployment without prospective validation on independently collected data (Discussion lines 713--717, Conclusion lines 749--750).

\change{Limitation item~6 at line 737; Discussion ``Honest performance assessment'' lines 713--717; Conclusion lines 749--750 (``neither is statistically superior to gradient boosting'').}

\bigskip
\hrule
\bigskip

\section*{Summary of All Changes Relevant to Reviewer PtZf}

\begin{center}
\small
\begin{tabular}{@{}p{3.8cm}p{7.7cm}@{}}
\toprule
\textbf{Concern} & \textbf{Revision (lines in \texttt{paper.tex})} \\
\midrule
Inconsistent reporting (W1) & \S4.3 L650--654 (reconciliation); Abstract L65; Table~1 L421--438 \\
Limited robustness (W2) & \S4.3 L593--599 (6-seed); Table~1 L432; Figure~4 L605; $p$-values L441, L507, L595; L598 (ensembling); \S5 L713--717 \\
Weak validation protocol (W3) & \S3.5 L314--318; L310; Limitation~4 L735 \\
Limited E3 generalization (W4) & Table~5 L625; Figure~5 L644; Root cause L612--620; Warning L619; Paradox L608--611; Abstract L65; Heading L608; Conclusion L749 \\
No external validation (W5) & Limitation~6 L737; \S5 L713--717; Conclusion L749--750 \\
\midrule
\multicolumn{2}{@{}l}{\textit{New analyses added in revision:}} \\
\midrule
6-seed extended validation & Table~1 L432; \S4.3 L593--599; Figure~4 L605 \\
Calibration analysis (ECE = 0.029) & \S4.6 L674--700; Table~6 L683 \\
Per-E3 breakdown & Table~5 L625; Figure~5 L644 \\
VHL root cause analysis & \S4.3 L612--620 \\
BRD4 case study & \S4.5 L655--672 \\
Lysine paradox discussion & \S5 L719--726; Figure~7 L729 \\
Statistical non-significance & L441, L507, L595 \\
\bottomrule
\end{tabular}
\end{center}

\end{document}
